\section{Introducción}
\label{sec:introduccion}

% ── A. Planteamiento general del tema ────────────────────────────────────────
\subsection{Planteamiento General del Tema}
\label{sec:intro_planteamiento}

La interacción entre las personas y los sistemas de atención al cliente atraviesa una
transformación sin precedentes. El sector de call centers está valuado en más de
400~mil millones de dólares, con una proyección de crecimiento a 496~mil millones para
2027 \cite{zoom2025}. En paralelo, la penetración de la inteligencia artificial
conversacional en estos entornos se acelera: el uso de \textit{speech analytics}
creció del 28\,\% al 37{,}5\,\% entre 2022 y 2023, evidenciando una demanda industrial
sostenida por herramientas de análisis automático de la voz \cite{zoom2025}. Este
contexto coloca el procesamiento afectivo del habla en el centro de una agenda de
innovación con impacto económico, operacional y humano real.

Sin embargo, la dimensión humana del problema ha sido subatendida. El 88\,\% de los
profesionales del sector reconoce el burnout de los agentes como uno de los mayores
desafíos de la industria, con tasas de rotación anual que oscilan entre el 38\,\% y el
52\,\% \cite{sqm2023, hiringbranch2024}. La causa estructural es clara: los mecanismos
de supervisión vigentes son reactivos y de cobertura limitada, incapaces de detectar
en tiempo oportuno los momentos de escalamiento emocional que preceden al agotamiento
crónico. El bienestar laboral, entendido como la capacidad de las organizaciones para
identificar y gestionar la carga emocional de sus agentes, emerge como una prioridad
que la tecnología tiene la obligación de atender.

En este escenario, la detección automática de picos emocionales en la voz constituye
un problema técnicamente maduro para ser abordado con métodos de aprendizaje automático
clásico combinados con análisis paralingüístico. Las características acústicas de la
voz ---frecuencia fundamental, intensidad, jitter, shimmer y coeficientes espectrales---
son correlatos fisiológicamente fundamentados de los estados afectivos del hablante;
y los algoritmos de detección de anomalías, aplicados sobre una línea base emocional
individualizada, ofrecen un mecanismo interpretable para señalar los momentos críticos
de cada llamada. Este trabajo explora si esa combinación es viable y con qué nivel de
rendimiento en el dominio específico del español de Colombia.

% ── B. Contextualización del problema ────────────────────────────────────────
\subsection{Contextualización del Problema}
\label{sec:intro_contexto}

Los sistemas actuales de gestión de calidad en call centers dependen de la escucha
manual por parte del supervisor, lo que limita la cobertura efectiva al 1--5\,\% del
total de interacciones \cite{zoom2025}; los sistemas de análisis de sentimiento basados
en texto, por su parte, son insensibles a las señales paralingüísticas que revelan
estados emocionales encubiertos como la frustración sutil o la ansiedad contenida
\cite{jha2022, feng2023}. Los detalles del problema, sus causas estructurales y sus
consecuencias organizacionales y humanas se desarrollan en la sección Área Problemática.

% ── C. Objetivo del estudio ───────────────────────────────────────────────────
\subsection{Objetivo del Estudio}
\label{sec:intro_objetivos}

\textbf{Objetivo General:} Desarrollar un sistema automatizado para detectar picos
emocionales en las llamadas de un call center, basado en el análisis de características
paralingüísticas del habla.

Para alcanzar este objetivo, se plantean cinco objetivos específicos: (1)~extraer y
caracterizar las features acústicas más relevantes ---pitch, intensidad, jitter,
shimmer, MFCCs y tasa de habla--- mediante Librosa, Parselmouth y OpenSMILE/eGeMAPS;
(2)~evaluar y comparar tres métodos para calcular líneas base acústicas individualizadas
por hablante (inicio de llamada, mediana global y ventana deslizante adaptativa);
(3)~comparar cuatro algoritmos de detección de anomalías ---Z-score estadístico,
Isolation Forest, One-Class SVM y Random Forest--- para identificar picos emocionales
sobre la línea base calculada; (4)~desarrollar un clasificador binario de valencia
emocional (positiva / negativa) que opere sobre los segmentos detectados como picos;
y (5)~validar el sistema completo sobre llamadas reales de call center en español con
anotaciones de expertos, reportando Precision, Recall y F1-score como métricas de
aceptación.

% ── D. Justificación resumida ─────────────────────────────────────────────────
\subsection{Justificación Resumida}
\label{sec:intro_justificacion}

El estudio se justifica por la convergencia de tres factores: la brecha de cobertura
industrial ---solo el 31\,\% de los contact centers monitorea activamente las emociones
del cliente \cite{zoom2025}---, la madurez tecnológica del SER ---con precisiones del
75--85\,\% en habla espontánea \cite{abdelhamid2022, wagner2023} que superan la línea
de viabilidad industrial--- y el costo humano del problema ---tasas de burnout
superiores al 63\,\% y rotaciones anuales de hasta el 52\,\% \cite{sqm2023,
hiringbranch2024}--- que hacen urgente una solución preventiva. La justificación
completa, incluyendo los beneficios operacionales, de bienestar laboral y académicos,
se presenta en la sección Justificación.

% ── E. Alcance ────────────────────────────────────────────────────────────────
\subsection{Alcance}
\label{sec:intro_alcance}

El sistema desarrollado en este proyecto \textbf{cubre}:
\begin{itemize}
    \item Llamadas telefónicas de call center en \textbf{español (Colombia)}.
    \item Dos canales de audio separados mediante diarización automática: agente y
          cliente.
    \item Detección de \textbf{picos de alta activación emocional}, definidos como
          desviaciones estadísticas significativas respecto a la línea base individual
          del hablante.
    \item \textbf{Clasificación binaria de valencia}: positiva (alegría, entusiasmo,
          satisfacción) o negativa (ira, frustración, ansiedad).
\end{itemize}

El sistema \textbf{no cubre}:
\begin{itemize}
    \item Clasificación de emociones específicas (ira, miedo, alegría, tristeza, etc.);
          solo se distingue positivo de negativo.
    \item Análisis multimodal: no se usa video, texto transcrito ni señales
          biométricas adicionales.
    \item Operación en tiempo real durante la llamada activa; el sistema es de análisis
          post-llamada (\textit{post-call analytics}).
    \item Análisis de la respuesta del agente ante el pico detectado.
    \item Generalización a dominios distintos del call center ni a idiomas distintos
          del español.
\end{itemize}

% ── F. Limitaciones ───────────────────────────────────────────────────────────
\subsection{Limitaciones}
\label{sec:intro_limitaciones}

Se reconocen de antemano las siguientes limitaciones del estudio:

\begin{itemize}
    \item \textbf{Calidad de audio:} las llamadas telefónicas se digitalizan a 8~kHz
          con códecs de compresión con pérdida, lo que reduce la resolución espectral
          disponible para la extracción de features de alta frecuencia.
    \item \textbf{Disponibilidad de corpus anotados en español:} los corpus de SER
          de referencia (RAVDESS, IEMOCAP) están en inglés y con habla actuada; el
          corpus de evaluación propio será de tamaño limitado (mínimo 30 llamadas).
    \item \textbf{Hardware:} el Dell Optiplex 9020 disponible restringe el uso de
          modelos de aprendizaje profundo de gran escala; la elección de métodos
          clásicos es en parte una adaptación a esta restricción.
    \item \textbf{Ventana temporal:} el proyecto se desarrolla en 15 semanas de
          práctica profesional, lo que limita la escala de los experimentos y la
          cantidad de iteraciones de ajuste.
    \item \textbf{Generalización fuera del dominio:} los resultados de rendimiento
          serán válidos para el dominio de call center en español colombiano; su
          transferibilidad a otros idiomas, acentos o sectores no está garantizada
          y no se evalúa en este trabajo.
\end{itemize}

% ── G. Estructura del documento ───────────────────────────────────────────────
\subsection{Estructura del Documento}
\label{sec:intro_estructura}

El resto del documento se organiza de la siguiente manera. La sección Referente
Teórico presenta los fundamentos del reconocimiento emocional en voz, las teorías
de las emociones, el modelo fuente-filtro, las características paralingüísticas
relevantes y la revisión de antecedentes. La sección Cronograma de Objetivos y la
sección Cronograma de Trabajo detallan las actividades y fases de implementación del
proyecto. La sección Metodología describe el diseño experimental, la población y
muestra, las herramientas, el procedimiento y el plan de validación con anotaciones
de expertos. La sección Resultados Esperados establece las métricas de aceptación y
los entregables técnicos. Finalmente, la sección Área Problemática profundiza en las
causas y consecuencias del problema, y la sección Justificación detalla los beneficios
esperados del estudio.
